\subsection{范畴表示论的基础}

\ssszh 有了这些一般性工具，我们来描述范畴表示论中的一些基本操作，并借此重新得到第 \ref{s:bbloc} 节中遇到的定理 \ref{t:loc+equiv} 与定理 \ref{t:lochchmod+eq}。

回忆：对代数群 $K$，有与之相伴的范畴表示的 $\infty$-范畴
%
\[
   \Dmod(K)\mod.
\]
%
表示论中别处也一样，$K$ 上的取逆诱导了 $\Dmod(K)$ 与其反向幺半 dg 范畴之间的同构；特别地，可以在左模与右模之间自由切换。

\ssszh 有一族基本的表示可以如下几何地构造出来。

给定带 $K$ 作用的簇 $X$，得到 $\Dmod(K)$ 在 $\Dmod(X)$ 上的典范作用。具体地说，底层的二元乘法由
%
\[
   \Dmod(K) \otimes \Dmod(X) \simeq \Dmod(K \times X) \xrightarrow{a_*}
   \Dmod(X)
\]
%
给出，与第 \ref{ex:actspace} 小节中函数论的情形完全一样。给定 $K$-簇的映射 $f: X \rightarrow X'$，沿 $f$ 的 $*$-推前与 $!$-拉回函子都带有典范的 $K$-等变数据。

\ssszh 特别地，取 $X = \on{pt}$，就得到 $K$ 在
%
\[
    \on{Vect} \simeq \Dmod(\on{pt})
\]
%
上的平凡表示。具体地说，若用 $\pi: K \rightarrow \on{pt}$ 记投影，则作用映射形如
%
\[
    \Dmod(K) \otimes \Vect \rightarrow \Vect, \quad \quad \sM \otimes V \mapsto
    \pi_*(\sM) \underset k \otimes  V.
\]
%

于是对 $K$ 的任意表示 $\sC$，可以作它的不变量与余不变量
%
\[
 \sC^K := \Hom_{\Dmod(K)\mod}(\on{Vect}, \sC) \quad \text{与} \quad \sC_K :=
 \on{Vect} \underset {\Dmod(K)} \otimes \sC.
\]
%
正如有限群通常表示的情形，二者典范地等同，论证如下。

由 $K$ 的光滑性，有 $\Dmod(K)$-等变的伴随对
%
\[
     \pi^![-2 \dim K]: \Dmod(K) \rightleftarrows \Vect: \pi_*,
\]
%
对它作用 $\Hom_{\Dmod(K)\mod}( -, \sC)$，便诱导出伴随对
%
\[
    \on{Oblv}: \sC^K \rightleftarrows \sC: \on{Av}_*.
\]
%
正如记号所暗示的，人们把 $\on{Oblv}$ 想成遗忘函子，把 $\on{Av}_*$ 想成平均函子。出于形式上的理由，平均函子经过余不变量分解：
%
\[
 \on{Av}_*: \sC_K \rightarrow \sC^K.
\]
%
根据 Beraldo 与 Gaitsgory 的一个很有用的定理 \cite{Bearaldo}、\cite{whitlocglob}，它是等价。

下面给出（余）不变量的一些基本例子。

\begin{ex}给定 $\Dmod(K)$ 的表示 $\sC$ 与 $\sD$，有典范等价
%
\[
       \Hom_{\Dmod(K)\mod}(\sC, \sD) \simeq \udHom(\sC, \sD)^K,
\]
%
其中的不变量关于 $\Dmod(K)$ 在映射空间上的对角作用来取。
\end{ex}

\begin{ex} \label{e:dmodquotstack} 对 $K$-簇 $X$，$\Dmod(X)$ 中 $K$-等变的对象典范地等同于商叠 $X/K$ 上 D-模的范畴，即
%
\[
     \Dmod(X)^K \simeq \Dmod(X/K).
\]
%
直白地说，D-模满足光滑下降，而图 \eqref{e:hombar} 正是沿覆盖 $X
\rightarrow X/K$ 计算 $X/K$ 上 D-模的「$\check{\on{C}}$ech 复形」。
\end{ex}

\begin{ex} 在 $\Dmod(K)$ 对其李代数表示范畴 $\fk\mod$ 的自然作用下，不变量等同于群本身的代数表示，即
%
\[
      \fk\mod^K \simeq K\mod.
\]
%
更一般地，给定同态 $K \rightarrow K'$，范畴
%
\[
 \fk'\mod^K
\]
%
典范地等同于对 $(\fk', K)$ 的 Harish-Chandra 模范畴。
\end{ex}

\begin{ex} 还有另一类例子，它们同时给出 $\fk\mod$ 上作用的另一种视角。设 $A$ 是带 $K$ 的自同构作用的 $k$-代数。此时得到 $\QCoh(K)$ 在 $A\mod$ 上的作用。

在单位元处求导，$K$ 在 $A$ 上的作用给出 $\fk$ 通过导子的作用 $\alpha$。再设这个作用有一个平凡化，即李代数的同态 $$\tau: \fk \rightarrow A,$$使得 $\alpha$ 与经 $\tau$ 给出的伴随作用一致。此时得到 $\Dmod(K)$ 在 $A\mod$ 上的作用。 \end{ex}

\begin{re}\label{r:invariantstwice}这里记录一个与通常表示情形基本不同之处。取不变量这一操作并非幂等，即平凡映射
%
\[
   \on{Oblv}_{\sC^K}: (\sC^K)^K \rightarrow \sC^K
\]
%
通常不是同构。取而代之的是典范等价
%
\[
   (\sC^K)^K \simeq (\sC^K) \otimes \Dmod(\on{pt/K}).
\]
%

之所以如此，是因为除非 $K$ 可缩，对任意（非零）表示 $\sD$，遗忘映射 $\on{Oblv}_{\sD}: \sD^K \rightarrow \sD$ 都不会是全忠实的，即 $\sD^K$ 与 $\sD$ 中的 Ext 并不一致。确实，在最普遍的 $\sD = \Dmod(K)$ 情形，有
%
\[
\Dmod(K)^K \simeq \Dmod(K / K) \simeq \Vect,
\]
%
其生成元是常值层。特别地，它在 $\Dmod(K)$ 中的自 Ext 是 $\pi_*(\underline{k})$，即 $K$ 的上同调；而它在 $\Vect$ 中的自 Ext 则是平凡的。

一般地，复合
%
\[
    \sD^K \xrightarrow{\on{Oblv}_{\sD}} \sD \xrightarrow{\on{Av}_*} \sD^K
    \xrightarrow{\on{Oblv}_{\sD}} \sD
\]
%
典范地等同于 $\on{Oblv}_{\sD} \otimes \pi_*(\underline{k})$，也就是把 $K$-等变对象底层的对象与 $K$ 的上同调作张量。读者不妨借助函数-层对应，把它与通常表示的平均过程中出现的有限群的基数因子作比较。

但要注意，若 $K$ 幂幺，则 $\on{Oblv}$ 是全忠实的，因为此时 $K$ 可缩，从而 $\pi_*(\underline{k}) \simeq k$。 \end{re}

\ssszh 讨论完不变量函子，我们来说明如何得到定理 \ref{t:loc+equiv} 与定理 \ref{t:lochchmod+eq}；先从后者开始。

给定代数群的映射 $H \rightarrow K$，有诱导函子与限制函子
%
\[
  \on{ind}_H^K: \Dmod(H)\mod \rightarrow \Dmod(K)\mod,
\]
%
%
\[
  \on{res}_K^H: \Dmod(K)\mod \rightarrow \Dmod(H)\mod.
\]
%
直白地说，$\on{res}_K^H$ 沿幺半函子
%
\[
         \Dmod(H) \rightarrow \Dmod(K)
\]
%
限制 $\Dmod(K)$ 在 dg 范畴 $\sC$ 上的作用，而 $\on{ind}_H^K$ 则把带 $\Dmod(H)$ 作用的 dg 范畴 $\sD$ 张量上去，即
%
\[
      \sD \mapsto \Dmod(K) \underset{ \Dmod(H) } \otimes \sD.
\]
%
正如有限群的 Frobenius 互反性，诱导与限制典范地互为左伴随，也互为右伴随。

\label{ss:invariantscosets} 特别地，得到 $\Dmod(K)$ 表示上的 $H$-不变量函子，即
%
\begin{align*}
     \big(\on{Res}_K^H(-)\big)^H &\simeq   \Hom_{\Dmod(H)\mod}( \on{Vect},
     \on{res}_K^H-).
\intertext{由伴随性，可以把它改写成 }
&\simeq  \Hom_{\Dmod(K)\mod}( \on{ind}_H^K \on{Vect},
-) \\ & \simeq \Hom_{\Dmod(K)\mod}( \Dmod(K) \underset{\Dmod(H)}{\otimes} \Vect
,
-)
\intertext{由 $H$ 的不变量与余不变量的等同、以及例 \ref{e:dmodquotstack}，可以认出这是}
& \simeq \Hom_{\Dmod(K)\mod}( \Dmod(K/H)
,
-).
\end{align*}

结论是：任意 $\Dmod(K)$ 表示的 $H$-不变量，经前复合而带有
%
\[
\Hom_{\Dmod(K)\mod}( \Dmod(K/H), \Dmod(K/H)) \simeq \Dmod(H \backslash K / H)
\]
%
的典范作用，也就是我们在第 \ref{s:bbloc} 节末尾遇到的 Beilinson--Drinfeld 的 Hecke 代数作用。

\ssszh 最后，对光滑的 $K$-簇 $X$，我们简述如何得到 $\Dmod(K)$-等变的伴随对
%
\[
    \on{Loc}: \fk\mod \rightleftarrows \Dmod(X): \Gamma(X, -).
\]
%
首先，从 $\OO_K$-模诱导到 $\sD_K$-模给出幺半函子
%
\[
      \QCoh(K) \rightarrow \Dmod(K).
\]
%
类似地，有 $\QCoh(K)$ 在 $\Dmod(X)$ 上的典范作用，使得遗忘函子 $\Dmod(X) \rightarrow \QCoh(X)$ 带有典范的 $\QCoh(K)$-等变数据。特别地，整体截面函子，即复合
%
\[
    \Dmod(X) \rightarrow \QCoh(X) \xrightarrow{\pi_*} \QCoh(\pt) \simeq \Vect,
\]
%
带有典范的 $\QCoh(K)$-等变数据。

然而，对任意 $\Dmod(K)$ 的表示 $\sC$，由伴随性有
%
\begin{align*}
\Hom_{\QCoh(K)\mod}( \on{Res} \sC, \Vect)
&\simeq \Hom_{\Dmod(K)\mod} \big(\sC,\\
&\qquad \Hom_{\QCoh(K)\mod}(\Dmod(K), \Vect)\big) \\
&\simeq \Hom_{\Dmod(K)\mod}(\sC, \fk\mod),
\end{align*}
%
其中最后的等价用了 $\fk\mod$ 与 $K$ 上弱 $K$-等变 D-模的等同，参见第 \ref{ss:weakinv} 小节。这就给出整体截面函子所需的 $\Dmod(K)$-等变性。至于局部化，可以利用 Gaitsgory \cite{whitlocglob} 的一个漂亮事实：$\Dmod(K)\mod$-等变函子的任意伴随都继承相容的等变数据。
